\begin{figure}[H]
\begin{multicols}{2}
    \scriptsize
    \begin{lstlisting}
\ COPYRIGHT 1991 JULIAN V. NOBLE
TASK INTEGRAL
FIND CP@L 0= ? ( FLOAD COMPLEX )
\ define data-type tokens if not already
FIND REAL*4 0= ?(((
    0 CONSTANT REAL*4
    1 CONSTANT REAL*8
    2 CONSTANT COMPLEX
    3 CONSTANT DCOMPLEX )))

FIND 1ARRAY 0= ?( FLOAD MATRIX.HSF )
\ function usage
: USE( [OOMPILE] ' CFA ; IMMEDIATE

\ BEHEADing starts here
0 VAR N

: inc.N N 1 + IS N ;
: dec.N N 2 - IS N ;

0 VAR type

\ define "stack"
20 LONG REAL*8   1ARRAY X{
20 LONG REAL*4   1ARRAY E{
20 LONG DCOMPLEX 1ARRAY F{
20 LONG DCOMPLEX 1ARRAY I{

2 DCOMPIEXSCALARS old.I final.I
: )imegral ( n-- ) \ trapezoidal rule
    X{ OVER    } G@L
    X{ OVER 1- } G@L
    F- F2/
    F{ OVER    } G@L
    F{ OVER 1- } G@L
    type2 AND
    IF   X+ FROT X*F
    ELSE F+ F* THEN
    I{ SWAP 1- } G!L ;
0 VAR f.name
: f(x) fname EXECUTE ;

: INITIALIZE
    IS type   \ store type
    type F{ !
    type I{ ! \ set types for function
    type ' old.I   !
    type ' final.I ! \ and integral(s)
    type 1 AND X{  !
            \ set type for indep. var.
    E{ 0 } G!L \ store error
    X{ 1 } G!L \ store B
    X{ 0 } G!L \ store A
    IS f.name  \ ! cfa of f(x)
    X{ 0 } G@L f(x) F{ 0 } G!L
    X{ 1 } G@L f(x) F{  1} G!L
    1 IS N
    N )integral
    type 2 AND IF F=0 THEN
    F=0 final.I G!L
    FINIT ;

: E/2 E{ N 1- } G@L F2/ E{ N 1- } G!L ;

: }move.down     ( adr n--)
    } \#BYTES >R ( --seg off )
    DDUP R@ +
                 ( --s.seg S.off d.seg d.off )
    R> CMOVEL ;

: MOVEDOWN
    E{ N 1- }move.down
    X{ N    }move.down
    F{ N    }mova.dcmn ;

: new.X ( 87:--x' )
    X{ N } G@L  X{ N 1- } G@L
    F+ F2/ FDUP X{ N    } G!L ;

\ cont'd. ...
    \end{lstlisting}
\end{multicols}
\end{figure}

\begin{figure}[H]
\begin{multicols}{2}
    \scriptsize
    \begin{lstlisting}
\ INTEGRAL cont'd
: GF.   1 > IF FSWAP E. THEN E. ;
: F@.   DUP>R G@L R> GF. ;
: .STACKS CR ." N"
     8 CTAB  ." X"
    19 CTAB  ." Re[F(X)]"
    31 CTAB  ." Im[F(X)]"
    45 CTAB  ." Re[I]"
    57 CTAB  ." Im[I]"
    71 CTAB  ." E"
    N 2 + 0 DO CR I .
         3 CTAB X{ I } F@.
        16 CTAB F{ I } F@.
        42 CTAB I{ I } F@.
        65 CTAB E{ I } F@.
    LOOP
    CR 5 SPACES ." old.I ="     old.I F@.
       5 SPACES ." final.I =" final.I F@. CR ;
CASE: <DEBUG> NEXT .STACKS ;CASE
0 VAR (DEBUG)
: DEBUG-ON  1 IS (DEBUG) 5 #PLACES
: DEBUG-OFF 0 IS (DEBUG) 7 #PLACES
: DEBUG (DEBUG) <DEBUG> ;

: SUBDIVIDE
    N 19 > ABORT" Too many subdivisions!"
    E/2 MOVE.DOWN
    I{ N 1- } DROP old.I #BYTES CMOVEL
        new.X f(x) F{ N } G!L
    N )integral N 1 + )integral ;

: CONVERGED? ( 87:--I[N]+I'[N-1]-I[N-1]:--f )
    I{ N } G@L I{ N 1- } G@L old.I G@L
    type 2 AND
    IF   CP- CP+ CPDUP CPABS
    ELSE  F-  F+  FDUP  FABS
    THEN
    E{ N 1- } G@L F2* F< ;

CASE: g*6 CP*F F* ;CASE
4 S->F 3 S->F F/ FCONSTANT F=4/3

: INTERPOLATE ( 87:I[N]=I'[N-1]-I[N-1]-- )
    F=4/3 type 2/ g*f
    old.I G@L final.I G@L
    type 2 AND
    IF   CP+ CP+
    ELSE  F+  F+ THEN
    final.I G!L ;
\ BEHEADing ends here

: )INTEGRAL ( 97:A B ERR--I[A,B] )
    INITIALIZE
    BEGIN N 0>
    WHILE   SUBDIVIDE   DEBUG
        CONVERGED?     inc.N
        IF INTERPOLATE dec.N
        ELSE type 2 AND
            IF   FDROP
            THEN FDROP
        THEN
    REPEAT final.I G@L ;
BEHEAD" N INTERPOLATE   \ optional
\ USE( F.name % A % B % E type )INTEGRAL
    \end{lstlisting}
\end{multicols}
\end{figure}